% GENERATED FILE. Edit canonical lessons, then run npm run generate:books.
% canonical-source-hash: fnv1a64:49f57a643ebc1a58
% canonical-lessons: GE-C15-sagen, GE-C15-gesagt, GE-C15-partizip, GE-C15-gemacht, GE-C15-gelernt, GE-C15-gewohnt, GE-C15-ge-praefix, GE-C15-partizip-practice

\chapter{The Wrapped Participle}
\label{ch:ge-partizip}
\hlchaptermodality{23}

\begin{chapteropening}
\textbf{By the end of this chapter:} \emph{I can turn any regular German verb I know into its past participle, and hear that the marking arrives at both ends of the word at once.}

The weak participle is a circumfix, so it is met as one concrete word first — gesagt — and only then as the recipe behind it. Three more verbs run through that recipe one at a time, the ge- gets its own lesson because it is the half English threw away, and the infinitive-to-participle table appears only at the end, where it recaps six lessons instead of introducing four words.
\end{chapteropening}

\section[sagen]{sagen — \textquotedblleft{}to say\textquotedblright{}}

\label{lesson:GE-C15-sagen}



\begin{quote}
One more weak verb before the past tense starts, and it is the verb the whole chapter will be built on.
\end{quote}

\subsection*{You'll want to know: sagen}
\begin{quote}
\textbf{sagen} — \textquotedblleft{}to say.\textquotedblright{}
\end{quote}

Two syllables, \emph{ZAH-gen}. It behaves like \emph{machen} and \emph{lernen} in every respect: chop the \emph{-en} and you have the stem \emph{sag-}, and onto that stem go the endings you already own — \emph{ich \textbf{sage}}, \emph{du \textbf{sagst}}.

Nothing new is happening in those endings. \emph{Sagen} is here so that when the past tense arrives in the next lesson, the only unfamiliar thing in it is the past.

\begin{sounds}
\begin{itemize}
\raggedright
  \item \texttt{long-a} — the \emph{a} of \emph{sagen} is long and open, the vowel of English \emph{father}, not the clipped \emph{a} of \emph{hat}.
  \item \texttt{ge-prefix} — the \emph{-gen} at the end is unstressed and swallowed, a soft \emph{gn} more than a full syllable. The same mumble will open every participle in this chapter, so it is worth getting comfortable now.
\end{itemize}
\end{sounds}

\begin{cousinweb}
\emph{Sagen} and English \emph{say} are the same word. Germanic *\emph{sagjaną} went to both, and then English quietly dropped the \emph{g} out of the middle:

\noindent
\begin{tabularx}{\linewidth}{@{}>{\raggedright\arraybackslash}X>{\raggedright\arraybackslash}X@{}}
\toprule
\textbf{German} & \textbf{English} \\
\midrule
\textbf{sag}en & \textbf{say} \\
\textbf{Tag} & \textbf{day} \\
\textbf{Weg} & \textbf{way} \\
\bottomrule
\end{tabularx}

Three pairs, one habit. Whenever a German \emph{g} sits between vowels or at the end of a short word, English is likely to have thinned it to a \emph{y} sound — so a German \emph{g} you cannot place is often an English \emph{y} in disguise.
\end{cousinweb}

\subsection*{Your turn}
Say these aloud:

\begin{itemize}
\raggedright
  \item \textquotedblleft{}sagen\textquotedblright{} — long \emph{a}, soft ending
  \item \textquotedblleft{}ich sage, du sagst, wir sagen\textquotedblright{}
  \item the three pairs — \textquotedblleft{}sagen–say, Tag–day, Weg–way\textquotedblright{}
\end{itemize}

\emph{Twice through:} \textquotedblleft{}sagen, sagen — to say.\textquotedblright{}

\begin{tcolorbox}[breakable,colback=teal!4,colframe=teal!35!black,title={Before you move on}]
Say \textquotedblleft{}to say.\textquotedblright{} (\textbf{Sagen}.) Say \textquotedblleft{}I say.\textquotedblright{} (\textbf{Ich sage}.) Is the \emph{a} long or short? (\textbf{Long} — as in \emph{father}.) Which English word is \emph{sagen}? (\textbf{Say}.) Next: what happens to this verb when the saying is already over.
\end{tcolorbox}
\section[gesagt]{gesagt — \textquotedblleft{}said\textquotedblright{}}

\label{lesson:GE-C15-gesagt}



\begin{quote}
You have \emph{sagen}. Now the form German uses when the saying is already finished — and it is longer at \textbf{both} ends than the word it came from.
\end{quote}

\subsection*{You'll want to know: gesagt}
\begin{quote}
\textbf{gesagt} — \textquotedblleft{}said.\textquotedblright{}
\end{quote}

Hear where it grew. \emph{Sag-} is still sitting in the middle, untouched. In front of it a \emph{ge-} has appeared; behind it the \emph{-en} has gone and a \emph{-t} has taken its place:

\begin{quote}
\emph{sagen} — \textbf{ge}-\emph{sag}-\textbf{t}
\end{quote}

That is the word English calls a \textbf{past participle}: the form that goes with \emph{have} in \emph{I have said}. German's has a hat on as well as boots.

\begin{sounds}
\begin{itemize}
\raggedright
  \item \texttt{ge-prefix} — unstressed, almost swallowed: \emph{guh}, never \emph{gay}. The stress stays on \emph{sagt}, exactly where it was in \emph{sagen}.
  \item \texttt{final-t} — the closing \emph{t} is crisp and fully released. The \emph{g} in front of it hardens to a \emph{k} on the way out, the same hardening you met when \emph{und} came out as \emph{unt}: say \emph{guh-ZAHKT}.
\end{itemize}
\end{sounds}

\subsection*{Your turn}
Say these aloud:

\begin{itemize}
\raggedright
  \item \textquotedblleft{}gesagt\textquotedblright{} — soft opening, hard ending
  \item the pair — \textquotedblleft{}sagen, gesagt\textquotedblright{}
  \item the three pieces — \textquotedblleft{}ge, sag, t\textquotedblright{}
\end{itemize}

\emph{Twice through:} \textquotedblleft{}sagen, gesagt. Sagen, gesagt.\textquotedblright{}

\begin{tcolorbox}[breakable,colback=teal!4,colframe=teal!35!black,title={Before you move on}]
Say \textquotedblleft{}said.\textquotedblright{} (\textbf{Gesagt}.) Which part of \emph{sagen} survives inside it? (\textbf{Sag-}, the stem.) What appeared in front? (\textbf{Ge-}.) What replaced the \emph{-en}? (\textbf{-t}.) Next: the same two additions, on any weak verb you like.
\end{tcolorbox}
\section[ge-…-t]{ge-…-t — the wrap that makes a participle}

\label{lesson:GE-C15-partizip}



\begin{quote}
\emph{Sagen}, \emph{gesagt}. One verb did that; the recipe behind it works on every regular verb you own.
\end{quote}

\begin{grammarlens}[title={Grammar Lens: ge- in front, -t behind}]
\begin{quote}
Take a regular verb. Strip the \emph{-en}. Put \textbf{ge-} in front of what is left and \textbf{-t} on the end.
\end{quote}

That is the whole operation for the regular verbs — the ones grammarians call \textbf{weak}, because they build their past out of endings rather than by changing the vowel inside.

The thing worth noticing is that the ending has a beginning. English marks its participle in one place: \emph{walk} becomes \emph{walk-\textbf{ed}}, and the front of the word sits untouched. German marks it in two, one piece of grammar arriving at both ends of the stem at once.

Grammar has a name for a piece that wraps like that: a \textbf{circumfix}. It is not a prefix plus a suffix that happen to co-occur — neither half means anything without the other. *\emph{Gesag} is not a word, and *\emph{sagt} on its own is \textquotedblleft{}he says.\textquotedblright{} Only both together make \textquotedblleft{}said.\textquotedblright{}
\end{grammarlens}

\subsection*{Your turn}
Say these aloud:

\begin{itemize}
\raggedright
  \item the recipe — \textquotedblleft{}ge in front, t behind\textquotedblright{}
  \item the wrap on its own — \textquotedblleft{}ge … t\textquotedblright{}
  \item \textquotedblleft{}sagen, gesagt\textquotedblright{} and hear both ends move
\end{itemize}

\emph{Twice through:} \textquotedblleft{}ge- in front, -t behind.\textquotedblright{}

\begin{tcolorbox}[breakable,colback=teal!4,colframe=teal!35!black,title={Before you move on}]
Which two pieces make a weak participle? (\textbf{Ge-} in front, \textbf{-t} behind.) What is a piece of grammar that wraps a word called? (A \textbf{circumfix}.) Where does English put its own participle marker? (\textbf{Only at the end} — \emph{walked}.) Build the participle of \emph{sagen}. (\textbf{Gesagt}.) Next: the same wrap on a verb you met long before this one.
\end{tcolorbox}
\section[gemacht]{gemacht — \textquotedblleft{}made,\textquotedblright{} \textquotedblleft{}done\textquotedblright{}}

\label{lesson:GE-C15-gemacht}



\begin{quote}
You know the recipe. Run it on \emph{machen} before reading on: strip the \emph{-en}, wrap what is left.
\end{quote}

\subsection*{You'll want to know: gemacht}
\begin{quote}
\textbf{gemacht} — \textquotedblleft{}made,\textquotedblright{} and also \textquotedblleft{}done.\textquotedblright{}
\end{quote}

\emph{Machen} gives up its \emph{-en}, leaving \emph{mach-}. The wrap goes round it and out comes \emph{gemacht}. No stem was harmed: the \emph{a} is the same \emph{a}, and the \emph{ch} is the same \emph{ch}.

German uses this one participle where English chooses between two different words. \emph{Ich habe das gemacht} is \textquotedblleft{}I made that\textquotedblright{} and \textquotedblleft{}I did that\textquotedblright{} alike; German never had to pick.

\begin{sounds}
\begin{itemize}
\raggedright
  \item \texttt{ch-ach} — after \emph{a}, \emph{o} or \emph{u}, German \emph{ch} is the scrape at the back of the throat, the sound in \emph{Bach} and in \emph{Nacht}. It sits in the middle of \emph{gemacht} untouched.
  \item \texttt{final-t} — and the \emph{t} on the end is crisp: \emph{guh-MAHKHT}, three consonants landing together at the finish.
\end{itemize}
\end{sounds}

\subsection*{Your turn}
Say these aloud:

\begin{itemize}
\raggedright
  \item \textquotedblleft{}machen, gemacht\textquotedblright{}
  \item both participles together — \textquotedblleft{}gesagt, gemacht\textquotedblright{}
  \item one word for two English ones — \textquotedblleft{}made, done — gemacht\textquotedblright{}
\end{itemize}

\emph{Twice through:} \textquotedblleft{}machen, gemacht.\textquotedblright{}

\begin{tcolorbox}[breakable,colback=teal!4,colframe=teal!35!black,title={Before you move on}]
Build the participle of \emph{machen}. (\textbf{Gemacht}.) Which two English words does it cover? (\textbf{Made} and \textbf{done}.) What happened to the \emph{ch} on the way? (\textbf{Nothing} — the wrap only touches the edges.) Next: the verb you have been doing to German itself.
\end{tcolorbox}
\section[gelernt]{gelernt — \textquotedblleft{}learned\textquotedblright{}}

\label{lesson:GE-C15-gelernt}



\begin{quote}
Third time through the recipe, and this one describes what you have been doing since page one.
\end{quote}

\subsection*{You'll want to know: gelernt}
\begin{quote}
\textbf{gelernt} — \textquotedblleft{}learned.\textquotedblright{}
\end{quote}

\emph{Lernen} loses its \emph{-en}, leaving \emph{lern-}; the wrap closes round it. The English cousin is doing the very same job with the very same stem — \emph{learn}, \emph{learned} — and differs only in where the marking sits. German puts a piece at each end; English puts both pieces at the back.

\begin{sounds}
\begin{itemize}
\raggedright
  \item \texttt{r-uvular-german} — the \emph{r} inside \emph{gelernt} is made at the back of the mouth and, before a consonant, is barely a consonant at all: it colours the vowel more than it rolls. \emph{Guh-LEHRNT}, with the \emph{r} almost swallowed.
  \item \texttt{final-t} — \emph{-nt} at the end, both consonants sounded, the \emph{t} released.
\end{itemize}
\end{sounds}

\subsection*{Your turn}
Say these aloud:

\begin{itemize}
\raggedright
  \item \textquotedblleft{}lernen, gelernt\textquotedblright{}
  \item the English twin — \textquotedblleft{}learn, learned; lernen, gelernt\textquotedblright{}
  \item all three so far — \textquotedblleft{}gesagt, gemacht, gelernt\textquotedblright{}
\end{itemize}

\emph{Twice through:} \textquotedblleft{}lernen, gelernt.\textquotedblright{}

\begin{tcolorbox}[breakable,colback=teal!4,colframe=teal!35!black,title={Before you move on}]
Build the participle of \emph{lernen}. (\textbf{Gelernt}.) Which English verb is its twin? (\textbf{Learn}.) Where does English put its marking? (\textbf{All of it at the end}.) Next: the fourth verb, and the last one before the recipe is yours.
\end{tcolorbox}
\section[gewohnt]{gewohnt — \textquotedblleft{}lived\textquotedblright{}}

\label{lesson:GE-C15-gewohnt}



\begin{quote}
Fourth and last of the four. Run the recipe on \emph{wohnen} yourself first.
\end{quote}

\subsection*{You'll want to know: gewohnt}
\begin{quote}
\textbf{gewohnt} — \textquotedblleft{}lived,\textquotedblright{} in the sense of dwelt somewhere.
\end{quote}

\emph{Wohn-} takes the wrap and nothing else changes. Four verbs, four participles, one recipe: it has not bent once.

\begin{sounds}
\begin{itemize}
\raggedright
  \item \texttt{w-is-v} — the \emph{w} buried in the middle is a German \emph{w}, which is English \emph{v}. \emph{Guh-VOHNT}.
  \item \texttt{final-t} — \emph{-nt} to close, exactly as in \emph{gelernt}.
\end{itemize}
\end{sounds}

\begin{cousinweb}
\emph{Wohnen} began by meaning something closer to \textquotedblleft{}to be accustomed\textquotedblright{} than \textquotedblleft{}to reside,\textquotedblright{} and \emph{gewohnt} still carries that older sense alongside the newer one: in modern German it is also the ordinary word for \emph{used to something}.

English kept the same word and only the old meaning:

\noindent
\begin{tabularx}{\linewidth}{@{}>{\raggedright\arraybackslash}X>{\raggedright\arraybackslash}X@{}}
\toprule
\textbf{word} & \textbf{meaning} \\
\midrule
German \textbf{gewohnt} & lived; accustomed \\
English \textbf{wont} & accustomed, as in \emph{he was wont to sing} \\
\bottomrule
\end{tabularx}

Two languages split one word between them — Germany took the address and England took the habit.
\end{cousinweb}

\subsection*{Your turn}
Say these aloud:

\begin{itemize}
\raggedright
  \item \textquotedblleft{}wohnen, gewohnt\textquotedblright{} — with a \emph{v} in the middle
  \item the fossil — \textquotedblleft{}gewohnt, wont\textquotedblright{}
  \item the whole set — \textquotedblleft{}gesagt, gemacht, gelernt, gewohnt\textquotedblright{}
\end{itemize}

\emph{Twice through:} \textquotedblleft{}wohnen, gewohnt.\textquotedblright{}

\begin{tcolorbox}[breakable,colback=teal!4,colframe=teal!35!black,title={Before you move on}]
Build the participle of \emph{wohnen}. (\textbf{Gewohnt}.) How is the \emph{w} sounded? (\textbf{Like English v}.) Which English word preserves its older sense? (\emph{\textbf{Wont}} — accustomed.) Next: where that \emph{ge-} came from, and the English word still hiding it.
\end{tcolorbox}
\section[ge-]{The ge- that English lost}

\label{lesson:GE-C15-ge-praefix}



\begin{quote}
Four participles, four \emph{ge-}s. That syllable is not decoration, and English used to have it too.
\end{quote}

\begin{cousinweb}
German \emph{ge-} comes from Germanic *\emph{ga-}, a prefix meaning \textquotedblleft{}\textbf{together, completely}.\textquotedblright{} Its job was to mark an action as \textbf{finished} — what grammarians call \emph{perfective} — which is exactly the work a past participle does. A finished saying is \emph{gesagt}.

English inherited the same prefix and spelled it \emph{y-}. Chaucer still wrote it. Then English let it go, and two fossils were left behind with the prefix welded on:

\noindent
\begin{tabularx}{\linewidth}{@{}>{\raggedright\arraybackslash}X>{\raggedright\arraybackslash}X@{}}
\toprule
\textbf{English word} & \textbf{what is hiding inside} \\
\midrule
\textbf{enough} & Old English \emph{ge}nōg — \textquotedblleft{}completely sufficient\textquotedblright{} \\
\emph{yclept} (archaic, \textquotedblleft{}named\textquotedblright{}) & \emph{ge}cleopod — \textquotedblleft{}called,\textquotedblright{} finished and done \\
\bottomrule
\end{tabularx}

Say the German wrap out loud — \emph{ge}-\emph{sag}-\emph{t} — and then the English fossil: \emph{en-ough}, from \emph{genōg}. It is the same syllable, worn down to a vowel and carried inside a word nobody thinks of as a participle at all.

So English once wrapped its participles the same way German still does. German simply never stopped.
\end{cousinweb}

\subsection*{Your turn}
Say these aloud:

\begin{itemize}
\raggedright
  \item what it meant — \textquotedblleft{}together, completely\textquotedblright{}
  \item the fossil — \textquotedblleft{}enough, from genōg\textquotedblright{}
  \item four finished actions — \textquotedblleft{}gesagt, gemacht, gelernt, gewohnt\textquotedblright{}
\end{itemize}

\emph{Twice through:} \textquotedblleft{}ge- means completely.\textquotedblright{}

\begin{tcolorbox}[breakable,colback=teal!4,colframe=teal!35!black,title={Before you move on}]
What did \emph{ge-} originally mean? (\textquotedblleft{}\textbf{Completely},\textquotedblright{} marking an action finished.) How did English spell the same prefix? (\emph{\textbf{Y-}}.) Which everyday English word still carries it? (\emph{\textbf{Enough}}, from \emph{genōg}.) And the archaic one? (\emph{\textbf{Yclept}}.) Next: the four participles side by side, and then the verb that puts them to work.
\end{tcolorbox}
\section[Practice]{Practice — four verbs, one wrap}

\label{lesson:GE-C15-partizip-practice}



\begin{quote}
No new words. Every participle here has had its own lesson, so this is the first place they are safe to put side by side — as a recap of work already done, not as a list to swallow.
\end{quote}

\begin{grammarlens}[title={Grammar Lens: the four, assembled}]
\noindent
\begin{tabularx}{\linewidth}{@{}>{\raggedright\arraybackslash}X>{\raggedright\arraybackslash}X@{}}
\toprule
\textbf{infinitive} & \textbf{participle} \\
\midrule
\emph{sagen} & \textbf{ge}sag\textbf{t} \\
\emph{machen} & \textbf{ge}mach\textbf{t} \\
\emph{lernen} & \textbf{ge}lern\textbf{t} \\
\emph{wohnen} & \textbf{ge}wohn\textbf{t} \\
\bottomrule
\end{tabularx}

Cover the right-hand side and rebuild it from the left. Then cover the left and work backwards: given \emph{gemacht}, what was the infinitive? The recipe runs both ways, which is what makes it worth having.
\end{grammarlens}

\subsection*{Your turn}
Say these aloud:

\begin{itemize}
\raggedright
  \item the four in order — \textquotedblleft{}gesagt, gemacht, gelernt, gewohnt\textquotedblright{}
  \item the four infinitives back — \textquotedblleft{}sagen, machen, lernen, wohnen\textquotedblright{}
  \item the recipe out loud — \textquotedblleft{}ge in front, t behind\textquotedblright{}
  \item the fossil — \textquotedblleft{}enough, from genōg\textquotedblright{}
\end{itemize}

\emph{Twice through:} \textquotedblleft{}gesagt, gemacht, gelernt, gewohnt.\textquotedblright{}

\begin{tcolorbox}[breakable,colback=teal!4,colframe=teal!35!black,title={Before you move on}]
Run the four participles without prompting. Which infinitive gives \emph{gemacht}? (\emph{\textbf{Machen}}.) What are the two pieces of the wrap? (\textbf{Ge-} and \textbf{-t}.) What did the front piece originally mean? (\textquotedblleft{}\textbf{Completely}.\textquotedblright{}) Next: the verb that carries all four of these into a sentence.
\end{tcolorbox}
